\chapter{The nutrient landscape}

\section{Learning objectives}
By the end of this chapter, the reader should be able to distinguish vitamins from minerals; explain the difference between intake, absorption, status, and requirement; interpret RDA, AI, EAR, UL, and DV cautiously; and identify clinical circumstances that make a dietary history more important than a supplement label.

\section{What vitamins and minerals are}
Vitamins are organic compounds required in small amounts for metabolism, tissue maintenance, gene regulation, blood formation, vision, immunity, and other processes. Minerals are elements. Some, such as calcium and potassium, are needed in comparatively large amounts; others, such as iodine and selenium, are needed in trace amounts. Both groups are essential, but neither supplies calories.

Essentiality means that an inadequate supply eventually produces a reproducible impairment of a biological function and that restoring the nutrient corrects or prevents the impairment. It does not mean that every person benefits from doses above physiological requirements. Nutrients act in networks: iron metabolism depends on inflammation and erythropoiesis; calcium handling depends on vitamin D, parathyroid hormone, kidney function, and bone remodeling; and folate and vitamin B12 share one-carbon metabolism. A single laboratory value therefore rarely tells the whole story.

\section{Classification and nomenclature}
The conventional list contains thirteen vitamins: A, C, D, E, K, and B1, B2, B3, B5, B6, B7, B9, and B12. A vitamin name may refer to a family of chemically related vitamers rather than one compound. For example, vitamin A activity includes retinoids and provitamin A carotenoids, while vitamin E includes tocopherols and tocotrienols. Some substances, including choline, are classified separately in many systems even though they share vitamin-like essentiality. Classification is useful for teaching but does not replace attention to chemical form, bioavailability, or dose. \cite{wikipediaVitamin}

\section{Historical perspective}
The clinical history of vitamins began with observations that particular foods prevented recognizable deficiency diseases: citrus fruits and scurvy, liver and night blindness, and less-polished grains and beriberi. In the early twentieth century, laboratory isolation and chemical characterization transformed these observations into a science of essential nutrients. The historical lesson is important: nutritional medicine developed from clinical observation, controlled dietary experiments, and biochemical investigation together. It also warns against assuming that a modern marketing claim has the evidentiary status of a proven deficiency treatment. \cite{wikipediaVitamin}

Vitamins are commonly divided into fat-soluble vitamins A, D, E, and K and water-soluble vitamin C and the B vitamins. Fat-soluble vitamins can be stored substantially in liver and body fat, so regular high-dose supplementation is more likely to accumulate. Water-soluble vitamins are not simply harmless in excess: vitamin B6, niacin, folic acid, and vitamin C can all cause problems at high intakes.

\section{Food, absorption, and status}
The amount eaten is not identical to the amount absorbed or used. Fiber, phytate, oxalate, cooking, fermentation, fat in a meal, gut health, medications, and the chemical form of a nutrient can change bioavailability. Blood tests can help in selected situations, but a result must be interpreted with symptoms, diet, medicines, and the relevant clinical context.

Nutritional epidemiology also requires care. Associations between a food pattern and a disease may reflect socioeconomic status, smoking, physical activity, healthcare access, or reverse causation. Randomized trials can test supplementation but may not reproduce the benefits of a whole-food pattern. In medical practice, the strongest conclusion is often specific: a nutrient prevents a deficiency-related outcome in a defined population, while evidence for generalized disease prevention remains uncertain.

\begin{table}[ht]
\centering
\caption{Useful terms}
\begin{tabular}{p{0.25\textwidth}p{0.65\textwidth}}
\toprule
Term & Meaning \\
\midrule
AI & Adequate Intake, used when evidence is insufficient to establish an RDA.\\
EAR & Estimated Average Requirement for a defined life-stage group.\\
RDA & Intake that meets the needs of nearly all healthy people in a group.\\
UL & Tolerable Upper Intake Level; the highest average daily intake unlikely to pose risk for most people.\\
DV & Daily Value used on food and supplement labels; it is a label reference, not a personalized target.\\
\bottomrule
\end{tabular}
\end{table}

\section{A sensible hierarchy}
Start with dietary variety, then address barriers such as low appetite, food insecurity, vegan eating, malabsorption, or a prescribed restriction. Consider a targeted supplement when there is a documented deficiency, a well-established life-stage need, or a clinician's recommendation. Treat symptoms such as fatigue, hair loss, cramps, or tingling as reasons for evaluation, not automatic proof of a deficiency.

\section{Assessment in clinical practice}
A useful nutritional assessment includes recent weight change, appetite, food access, food preparation, alcohol intake, gastrointestinal symptoms, chewing or swallowing difficulty, menstrual or other blood loss, pregnancy status, medications, supplements, and relevant diagnoses. Examination may look for pallor, glossitis, dermatitis, edema, muscle wasting, bone tenderness, or goiter, but these findings are neither sensitive nor specific. Laboratory testing should be selected to answer a clinical question rather than performed as an indiscriminate micronutrient panel.

\section{Key points}
\begin{itemize}
\item Requirements are population reference values; they are not automatic treatment targets.
\item A varied diet improves the probability of adequacy and supplies fiber, protein, and bioactive compounds together.
\item Deficiency and toxicity are clinical states with causes, manifestations, and management plans.
\item Supplement decisions should include dose, duration, interaction review, and reassessment.
\end{itemize}

\warning{The nutrient chapters summarize general evidence. They do not replace a clinician, especially during pregnancy, in children, in kidney or liver disease, or when taking prescription medicines.}
